\section{Estados y flujos}

\subsection{Descripcion}

Este documento define los estados funcionales del sistema y las transiciones validas entre ellos. Su foco es el dominio, no la implementacion tecnica de pantalla o endpoint.

\subsection{Alcance}

Aplica a usuarios, partidos, picks, resultados y torneo.

\subsection{Definiciones}

\begin{itemize}
  \item \textbf{Transicion}: cambio valido entre dos estados.
  \item \textbf{Estado terminal}: estado desde el cual no se espera retorno ordinario.
\end{itemize}

\subsection{Reglas}

Estados de usuario:

\begin{itemize}
  \item `PENDIENTE`: registrado, con pago aun no validado o en revision.
  \item `ACTIVO`: habilitado para participar.
  \item `RECHAZADO`: activacion denegada.
  \item `BLOQUEADO`: acceso revocado por accion administrativa.
\end{itemize}

Transiciones validas de usuario:

\begin{itemize}
  \item registro nuevo a `PENDIENTE`;
  \item `PENDIENTE` a `ACTIVO`;
  \item `PENDIENTE` a `RECHAZADO`;
  \item `ACTIVO` a `BLOQUEADO`.
\end{itemize}

Estados de partido a nivel funcional:

\begin{itemize}
  \item `PROGRAMADO`;
  \item `BLOQUEADO` para picks;
  \item `CERRADO`;
  \item `SUSPENDIDO`;
  \item `POSTERGADO`;
  \item `CANCELADO`;
  \item `CORREGIDO` como marca de resultado ajustado.
\end{itemize}

Transiciones validas de partido:

\begin{itemize}
  \item `PROGRAMADO` a `BLOQUEADO`;
  \item `BLOQUEADO` a `CERRADO`;
  \item `PROGRAMADO` o `BLOQUEADO` a `POSTERGADO`;
  \item `BLOQUEADO` o `CERRADO` a `SUSPENDIDO`;
  \item `PROGRAMADO`, `BLOQUEADO` o `SUSPENDIDO` a `CANCELADO`;
  \item `CERRADO` a `CORREGIDO`.
\end{itemize}

Estados de pick:

\begin{itemize}
  \item `ABIERTO`: editable;
  \item `BLOQUEADO`: no editable;
  \item \texttt{SIN\_PICK}: no existe prediccion guardada al momento del lock.
\end{itemize}

Estados de torneo:

\begin{itemize}
  \item `ACTIVO`;
  \item `ARCHIVADO`.
\end{itemize}

\subsection{Casos borde}

\begin{itemize}
  \item Un usuario rechazado no se reactiva sin una accion administrativa explicita y documentada.
  \item Un partido corregido conserva su historial previo mediante auditoria; la correccion no borra la trazabilidad.
  \item \texttt{SIN\_PICK} no es una prediccion; es la ausencia valida de una prediccion antes del cierre.
\end{itemize}

\subsection{Invariantes}

\begin{itemize}
  \item Un usuario no puede ser simultaneamente `ACTIVO` y `BLOQUEADO`.
  \item Un torneo archivado no vuelve a estado activo dentro del flujo ordinario.
  \item Un pick bloqueado no vuelve a estado abierto por flujo ordinario.
\end{itemize}

\subsection{Preguntas abiertas}

\begin{itemize}
  \item Confirmar si se necesita un estado funcional explicito para pago \texttt{EN\_REVISION} como entidad separada del estado de usuario.
\end{itemize}

\subsection{Decisiones pendientes}

\begin{itemize}
  \item `Decision pendiente` `No bloqueante`: formalizar los estados internos del comprobante de pago si el backend los modela como entidad propia.
\end{itemize}
